\chapter{Introduction}
\label{chap:intro}

\section{Background and Motivation}

Road traffic safety is among the foremost public health concerns in developing countries.
India accounts for a disproportionately large share of global road fatalities, with a
significant proportion of accidents occurring at uncontrolled or poorly managed intersections.
Roundabouts --- circular junctions that regulate traffic through yield-based priority rather
than traffic signals --- are structurally well-suited to reducing conflict severity; however,
the absence of signal enforcement in Indian roundabouts leads to persistent unsafe behaviours:
aggressive overtaking, tailgating, wrong-way entry, abrupt stopping within the circulatory
zone, and erratic lane weaving.

Manual monitoring of such intersections is resource-intensive, subjective, and temporally
limited.  Unmanned aerial vehicles (UAVs) equipped with downward-facing cameras offer a
compelling alternative: a stationary overhead drone captures the full spatial extent of an
intersection without the occlusion that afflicts ground-mounted sensors, and records a
continuous, repeatable traffic record amenable to offline automated analysis.

\section{Problem Statement}

Given traffic video captured from a drone at an unsignalized roundabout intersection, it is
difficult to automatically characterise traffic safety because raw video does not directly
provide structured information about which road users are present, where they are located,
how they move over time, how vehicles interact spatially, whether their behaviour violates
predefined safety conditions, or which traffic situations should be considered dangerous.

Addressing this gap requires three tightly coupled sub-problems to be solved in sequence.

\begin{enumerate}
    \item \textbf{Structured trajectory extraction:}
          Raw video pixels must be transformed into reliable, persistent per-vehicle
          trajectories.  Dense Indian intersection footage presents extreme challenges for
          standard detection and tracking pipelines: vehicles occupy as few as $15 \times 8$
          pixels at typical drone altitudes; fifty or more road users may be simultaneously
          present in a single frame; bounding boxes frequently overlap; and partial occlusion
          causes confidence drops that fragment track identities across time.

    \item \textbf{Meaningful spatial representation:}
          Trajectories expressed in image/pixel coordinates are not physically meaningful.
          Safety rules involve real-world distances (e.g., following distance in metres) and
          speeds (e.g., m/s).  Converting pixel coordinates to world coordinates requires a
          calibrated pixel-to-metre mapping grounded in actual physical measurements at the
          intersection.

    \item \textbf{Safety assessment and conflict prediction:}
          Metric trajectories must be interpreted through explicitly designed safety rules
          that identify dangerous behaviours.  The outputs of these rules, combined with
          manually annotated dangerous traffic events, must then be used to train a machine
          learning model capable of predicting dangerous vehicle interactions.
\end{enumerate}

This project develops an end-to-end pipeline that addresses all three sub-problems within a
unified system architecture tailored to the specific characteristics of Indian roundabout
traffic.

\section{Objectives}

\begin{enumerate}
    \item To develop a high-recall vehicle detection pipeline for small-object, high-density
          aerial imagery using SAHI tiling and a DINO-style RT-DETRv2-X detection backbone.
    \item To implement a robust multi-object tracker (ByteTrack) augmented with appearance-based
          track association that maintains persistent track identities across the video duration.
    \item To calibrate pixel coordinates to real-world metric coordinates using physical
          road measurements obtained at the actual intersection, enabling physically
          meaningful speed and distance computation.
    \item To design and implement a deterministic safety rule engine that evaluates five
          categories of traffic violation --- tailgating, wrong-way driving, erratic lane
          weaving, unsafe overtaking, and vehicle stoppage --- using geometric and temporal
          conditions derived from the metric trajectory data.
    \item To construct a training dataset from safety-rule outputs and manually annotated
          dangerous traffic events, and to train a calibrated machine learning model capable
          of assigning probabilistic danger scores to vehicle interaction pairs.
\end{enumerate}

\section{Scope of Work}

The project analyses a set of traffic video recordings of an Indian roundabout intersection
captured from a stationary overhead drone.  All processing is offline; real-time inference
is not a requirement.  The three pipeline components --- detection and tracking, safety rule
evaluation, and machine learning --- operate on these recordings but are designed to generalise
to similar overhead footage.

\section{Organisation of the Report}

Chapter~\ref{chap:lit_review} reviews related work.
Chapter~\ref{chap:proposed} presents the system architecture and data flow.
Chapter~\ref{chap:detection} details the detection, tracking, calibration, and trajectory
generation methodology.
Chapter~\ref{chap:safety} documents each safety rule.
Chapter~\ref{chap:ml} explains the dataset construction and machine learning model.
Chapter~\ref{chap:results} presents results.
Chapters~\ref{chap:discussion} and~\ref{chap:conclusion} discuss limitations and conclude.